\section{The tree selected-inverse schedule and its projected adjoint}
\label{app:theorem}

This appendix states the classical machinery used in
Section~\ref{sec:theorem}: the tree recurrence that makes the layer exact and
global, and the projected-Hessian argument for its first-order adjoint.

\subsection{Forward: a two-pass schedule on the elimination tree}

\begin{theorem}[Tree selected inverse]
\label{thm:tree-forward}
Let $\Gr$ be a tree rooted at $r$ and $A$ SPD with pattern $\Gr$. Order nodes
so that each precedes its parent. Then \textbf{(a) zero fill:} the block
$LDL\T$ factorization in this order has $L$ unit lower-triangular with the
pattern of $A$ and every pivot $\piv{v}$ SPD; \textbf{(b) collect
(leaves$\to$root):}
$\piv{v} = \Ablock{v}{v} - \sum_{c\in\children(v)} \Ablock{v}{c}\inv{\piv{c}}\Ablock{c}{v}$
and $\Gblock{r}{r} = \inv{\piv{r}}$; \textbf{(c) distribute
(root$\to$leaves):} with $\fac{v} = \Ablock{\parent(v)}{v}\inv{\piv{v}}$,
\[
  \Gblock{\parent(v)}{v} = -\,\Gblock{\parent(v)}{\parent(v)}\,\fac{v},
  \qquad
  \Gblock{v}{v} = \inv{\piv{v}} + \fac{v}\T \Gblock{\parent(v)}{\parent(v)} \fac{v}.
\]
This produces every block of $\selinv(A)$ exactly in
$\bigO((|V|+|E|)\,b^3)$ time and $\bigO((|V|+|E|)\,b^2)$ memory. Steps
(b)--(c) are the tree specialization of the classical Gaussian-BP
collect/distribute and Takahashi recurrences
\citep{weiss2001correctness,takahashi1973}.
\end{theorem}

\subsection{Why the derivative has a same-pattern reverse schedule}

\begin{theorem}[Projected selected-inverse derivative is self-adjoint]
\label{thm:p0}
On the SPD cone, under the full Frobenius inner product on symmetric
matrices, $\selinv_{\pat}(A) = P_{\pat}\nabla_A \log\det A$.
Consequently the Fr\'echet derivative
$D\,\selinv|_A = P_{\pat}\,L_A\,\iota_{\pat}$ (with
$L_A: X\mapsto -\inv{A}X\inv{A}$ the derivative of inversion and $P_{\pat}$
the on-pattern projection, and $\iota_{\pat}$ the inclusion of symmetric
on-pattern perturbations) is self-adjoint: it is the projected Hessian of a
scalar potential, and a gradient field has a symmetric Jacobian (symmetry of
second derivatives).
Hence, for the tree schedule above, the reverse-mode map
$\cotan{G}\mapsto\cotan{A}$ equals the forward derivative map and is realized
by reversing the schedule of Theorem~\ref{thm:tree-forward} at the same structural work
(cheap-gradient principle \citep{baur1983complexity}).
\end{theorem}

\begin{remark}[Self-adjoint up to transpose; the non-symmetric correction]
\label{rem:transpose}
For non-symmetric $A$,
$L_A^* : X\mapsto -A^{-\mathsf{T}}XA^{-\mathsf{T}} = L_{A\T}$, so the adjoint
is $P_{\pat}L_{A\T}\iota_{\pat}$: the projected inversion derivative
evaluated at $A\T$. It is not itself the selected inverse of $A\T$. The
non-symmetric (LU / Erisman--Tinney) kernels are therefore not literally
self-adjoint; the derivative is adjoint up to transpose, which is why
their backward propagates independent lower and upper factor cotangents (the
$L$ of $A$ is the $U$ of $A\T$).
\end{remark}

\begin{remark}[Why the naive off-pattern argument fails]
It is tempting to argue that $\cotan{A}|_{\pat}$ never references an
off-pattern block of $\inv{A}$. That is false:
$\cotan{A}_{vv} = -\sum_{(s,t)\in\pat} \Gblock{v}{s}\cotan{G}_{st}\Gblock{t}{v}$
involves off-pattern $\Gblock{v}{s}$. What is true is that the reverse
schedule never forms them: it propagates cotangents along tree edges, exactly
as the forward computes $\Gblock{v}{v}$ without forming the dense inverse.
The first-order reverse schedule closes on the pattern. This statement alone
does not establish pattern closure or the cost of arbitrary higher-order
derivatives.
\end{remark}

\section{Engine details: generality, stability boundary, and a drop-in}
\label{app:engine}

\paragraph{No-pivot stability boundary (non-symmetric).}
The non-symmetric kernels eliminate on a static pattern with no pivoting,
valid while pivots stay non-singular.
On the tested sweep this is at parity with a pivoted dense LU while the
system is block-diagonally dominant, with the accuracy floor collapsing below the
dominance boundary. The SPD experiment shows a small reconstructed-factor
residual on the finite tree/conditioning grid tested; it does not prove a
height- or $\kappa$-independent global backward-error theorem. The
non-symmetric diagnostic additionally exposes the no-pivot boundary. These
are numerical-analysis observations with no demonstrated learning payoff,
which is why they are not in the body. Data:
\href{paper/attainability/figures/data/nonsym_stability.csv}
{\texttt{nonsym\_stability.csv}} and
\href{paper/attainability/figures/data/stability.csv}
{\texttt{stability.csv}}; generator:
\href{paper/attainability/make_figures.py}{\texttt{make\_figures.py}}.

\paragraph{DeltaNet chunk inverse (a drop-in capability note).}
The differentiable triangular chunk inverse $(I-A)^{-1}$ drops, unchanged,
into a chunked delta-rule linear-attention layer
\citep{yang2024parallelizing}: the within-chunk delta rule \emph{is} that
triangular solve. It reproduces the token-by-token delta rule exactly, and
its analytic backward equals autograd through the whole multi-chunk layer.
Because the forward is identical by construction, this demonstrates
availability (the operator is a genuine drop-in), not a capability gain, and
so it is recorded here rather than as a result. A learned-capability payoff for
sequence models would be a learned-sparse graph-coupled mixer, which belongs
to the program of Section~\ref{sec:program} rather than to the present
evidence.

\section{Confirmatory evidence chronology}
\label{app:integrity}

At repository commit \evidencecommit{}, the matched-fit protocol was frozen
on 29 June 2026, before the first aggregate read of HOLDOUT seeds 1000--1029.
The evidentiary records are \url{archive/docs/CONFIRMATORY_RESULTS.md} and
\url{archive/results/CONFIRMATORY/matched_fit_taskA/}; only these support the
paper's confirmatory interpretation.

On 30 June, a later branch froze a bootstrap protocol and reused the same
HOLDOUT seeds. Its assertion that the seeds were untouched is false in
wall-clock chronology, so that use is not an independent confirmation and
its verdict and effect-size claims are excluded. Both protocol files retain
dated provenance notes.

The first-read \texttt{summary.json} records a top-level
\texttt{run\_tag=CONFIRMATORY}, the HOLDOUT seed list, and a confirmatory
output path, but two nested payloads say \texttt{EXPLORATORY}. In
\url{archive/bench/matched_fit_taskA.py}, those labels come from a module-level
constant; seed guarding and output routing use the confirmatory function
argument. This metadata defect does not change a metric and remains disclosed
in the frozen record.

The maze summary contains a second ambiguity. Its \texttt{floor\_sound} field
is \texttt{false}: at the in-distribution size, the $K=256$ floor arm's mean
evaluation-loss gap from exact is $2.70\times10^{-5}$, versus the
implementation's $3.07\times10^{-6}$ threshold. Results prose nevertheless
marked maze NC2 as passing because its mean absolute extrapolation gaps were
3.8--4.6 orders below the $K=7$ gaps, the narrower check implemented by
\texttt{confirmatory\_decision}. The preregistration required collapse in
``function/extrapolation distance'' and did not specify which of these
non-equivalent summaries controlled. This manuscript does not retroactively
choose the favorable one: maze NC2 is ambiguous, and no positive conclusion
depends on treating it as a pass.

A third defect is structural and was found in a 2026-07-25 re-audit of the
frozen adjudication code. The preregistration defines
$F=\max(2M_{\mathrm{null}},2\times10^{-6})$, where
$M_{\mathrm{null}}$ is the near-exact cell's own matched function distance,
and applies $F$ to both that cell's C2 collapse test and the rival
inductive-bias branch. At high difficulty, $K=32$ is the only arm reaching
the $60\%$ matchability bar and thus the only adjudicable equal-fit arm; it
also sets $M_{\mathrm{null}}$, so its distance is compared with twice itself.
In
\url{archive/bench/matched_fit_taskA.py}, the verdict is \textsc{survive} when
the matched distance's confidence interval (CI) lower bound exceeds $F$,
\textsc{collapse} when its mean is at most $F$, and \textsc{ambiguous}
otherwise. With $F=2\times$ that mean, \textsc{collapse} is forced and
\textsc{survive} is unreachable. C2's recorded pass is therefore
uninformative, and the inductive-bias branch could not have fired on these
data regardless of the measurements. Measured instead against the
preregistration's absolute optimization-noise scale $2\times10^{-6}$,
calibrated on development seeds where the largest silent-cell matched distance
was $1.46\times10^{-6}$, the two near-exact distances
($9.11\times10^{-4}$ at $\rho=0.98$ and $1.83\times10^{-4}$ at $\rho=0.99$)
are $455\times$ and $91\times$ larger, and the $\rho=0.98$ CI lower bound
($3.84\times10^{-4}$) also exceeds it. We report the frozen rule unchanged.
Thus, the evidence does not support an inductive-bias reading where the finite
arms cannot fit and leaves it untested where they can.
